% !TEX root =./ECE343_Practice_main.tex
%
% L1 practice set (Course Introduction): the spectrum and wavelength, how the
% course is assessed, and the circuit laws as limits of Maxwell's equations.

\fullwidth{\textbf{Learning Objectives 1.1 and 1.2: } I can articulate what I will
learn and how I will be assessed in this course. I can articulate practical
applications of electromagnetics.
\\[4pt]\normalfont\small Show your work. Write a \textbf{Documentation} line at the
end (\textbf{None} if you did not collaborate). Use $c = 3\times10^{8}$ m/s.}

\begin{questions}

%%%%%%%%%%%%%%%%%%%%%%%%%%%%%%%%%%%%%%%%%
\question \textbf{Wavelengths across the spectrum.}
\begin{parts}

	\begin{minipage}[t]{0.58\linewidth}
		\part Find the free-space wavelength of the $60\ \hz$ power grid, the GPS L1
		signal at $1575.42\ \mhz$, and an X-band radar at $10\ \ghz$.
		\begin{solution}[1.3in]
		\eq{
			\lambda_{60\ \text{Hz}} &= \dfrac{3\times10^{8}}{60} = 5\times10^{6}\ \m = 5000\ \km \\
			\lambda_\text{GPS} &= \dfrac{3\times10^{8}}{1.57542\times10^{9}} = 0.190\ \m \\
			\lambda_\text{X} &= \dfrac{3\times10^{8}}{10^{10}} = 0.03\ \m
		}
		\end{solution}
	\end{minipage}\hfill%
	\begin{minipage}[t]{0.37\linewidth}
		\raggedleft \vspace{12pt}
		$\lambda_{60\ \text{Hz}}$ = \ansbox[1.3in]{$5000\ \km$}

		\vspace{4pt}
		$\lambda_\text{GPS}$ = \ansbox[1.3in]{$19.0\ \text{cm}$}

		\vspace{4pt}
		$\lambda_\text{X}$ = \ansbox[1.3in]{$3\ \text{cm}$}
	\end{minipage}

	\part A $1\ \m$ connection carries each of the three signals. For which of them is
	the connection a significant fraction of a wavelength? Explain in one or two
	sentences why that matters.
	\begin{solution}[1in]
	At $60\ \hz$ the connection is $2\times10^{-7}$ of a wavelength, so it behaves as
	an ideal wire. At GPS and X-band frequencies it is about 5 and 33 wavelengths
	long, so the voltage differs from one end to the other at each instant, and
	circuit theory no longer describes it; Lesson 3 models it as a transmission
	line.
	\end{solution}

	\part Name one Air Force system that uses electromagnetic waves, and state one
	question about it that this course will help you answer.
	\begin{solution}[0.9in]
	Answers vary. Example: an aircraft radar. How much of the transmitter's power
	reaches the antenna through the feed cable, and how much reflects at a
	mismatched connection?
	\end{solution}

\end{parts}

\newpage
%%%%%%%%%%%%%%%%%%%%%%%%%%%%%%%%%%%%%%%%%
\question \textbf{The course, and where circuit theory comes from.}
\begin{parts}

	\part How is each learning objective assessed, and what can you do if you earn
	a 0 on an assessment?
	\begin{solution}[0.9in]
	Each objective has an assessment scored 1 (mastered) or 0 (not yet mastered).
	A 0 comes with feedback, and if you made a real attempt on time you may revise
	your work until you earn the 1.
	\end{solution}

	\part A transformer steps $120\ \text{V}$ down to $12\ \text{V}$. Which of
	Maxwell's equations explains how it works? Explain in one sentence.
	\begin{solution}[0.8in]
	Faraday's law: the changing magnetic flux in the core, produced by the primary
	current, induces a voltage around the secondary winding.
	\end{solution}

	\part Faraday's law in integral form is
	$\oint_C \vec{E}\cdot d\vec{\ell} = -\dfrac{d}{dt}\int_S \vec{B}\cdot d\vec{s}$.
	Explain how it reduces to Kirchhoff's voltage law, and under what condition.
	\begin{solution}[1.1in]
	Around a circuit loop, the line integral of $\vec{E}$ across each element is the
	voltage across it, so the left side is the sum of the voltages around the loop.
	When the magnetic flux through the loop changes negligibly, the right side is
	zero, and the sum of the voltages around the loop is zero.
	\end{solution}

	\part Kirchhoff's current law follows from conservation of charge. State the
	condition under which it holds, and give one situation in this course where it
	does not.
	\begin{solution}[1in]
	It holds when no charge accumulates at the node, so the current in equals the
	current out. On a transmission line at high frequency, charge builds up and
	drains along the conductors as the wave passes, so the current varies along the
	line; the line's capacitance accounts for it.
	\end{solution}

\end{parts}

\vspace{1.5em}
\noindent\textbf{Documentation:} \rule[-2pt]{0.6\linewidth}{0.4pt}

\end{questions}
